\documentclass[a4paper,10pt]{report}\input{../0.Préambule/Préambule_cours_prof}\begin{document}

\pagestyle{empty}
\newpage

\section*{Proportionnalité et pourcentage}
\addcontentsline{toc}{section}{Proportionnalité et pourcentage}

\noindent \textit{La propreté de la copie et de la rédaction sera \textbf{noté sur 2 points} }

\noindent \textit{L'usage de la calculatrice est \textbf{autorisé}.} \hfill \textit{Durée : 55 minutes}

\vspace{-3mm}
\paragraph{Exercice 1}: \hfill \textbf{1.5 points}

\noindent Le père de Lucas achète une tranche de $150$g de Roquefort Bio à $3$ euros,
\begin{enumerate}
\item Combien coûtent $100$g de ce fromage ? et $300g$ ?
\item Quel est le prix au kilogramme de ce fromage ?
\end{enumerate}

\vspace{-3mm}
\paragraph{Exercice 2}: \hfill \textbf{1.5 points}

\noindent Gabriella participe à une épreuve de course à pieds. Elle court à la même vitesse tout au long du parcours. Les résultats des temps de passage figurent dans la tableau suivant :

\medskip
\noindent \begin{minipage}{9cm}
\begin{enumerate}
\item Ce tableau est-il un tableau de proportionnalité ? Justifie
\item Si oui, calcule le coefficient de proportionnalité ?
\item Quelle distance Gabriella parcourt-elle en une heure ?
\end{enumerate}
\end{minipage}
\hspace{2mm}
\begin{minipage}{7.5cm}
\begin{center}
\renewcommand{\arraystretch}{1.5}
\begin{tabularx}{\linewidth}{|p{4cm}|*{4}{>{\centering \arraybackslash}X|}}
\hline
\cellcolor{lightgray!50}{
Kilomètres effectués} & $4,5$ & $6$ & $11$ & $14$\\\hline
\cellcolor{lightgray!50}{Temps en minutes} & $27$ & $36$ & $66$ & $84$\\\hline
\end{tabularx}
\end{center}
\end{minipage}

\vspace{-3mm}
\paragraph{Exercice 3}: \hfill \textbf{1 points}

\noindent Calcule les pourcentages suivants :
\begin{enumerate}
\begin{tabular}{m{4cm}m{4cm}m{4cm}m{4cm}}
\item $25\%$ de $12$ &
\item $75\%$ de $52$L &
\item $50\%$ de $438$m &
\item $60\%$ de $200$km \\
\end{tabular}
\end{enumerate}

\vspace{-3mm}
\paragraph{Exercice 4}: \hfill \textbf{1.5 points}

\noindent Au mois de janvier, le loyer des parents de François a augmenté de $4 \%$. Le loyer était de $800$ euros.
\begin{enumerate}
\item A combien s'élève l'augmentation ?
\item Quel est alors le montant du nouveau loyer ?
\end{enumerate}

\vspace{-3mm}
\paragraph{Exercice 5}:  \hfill \textbf{1 points}

Au cours du dernier trimestre, une usine d'électro-ménager a produit $1 400$
téléviseurs. Le service après-vente a noté des dysfonctionnements sur $119$ d'entre eux. Quel est le pourcentage d'électro-ménager qui dysfonctionne qui ont des lunettes ?

\vspace{-3mm}
\paragraph{Exercice 6}: \hfill \textbf{1 points}

Dans une classe, il y a $24$ élèves. Parmi ces élèves, $9$ ont des lunettes.
Quel est le pourcentage d'élèves qui n'ont pas de lunette ?

\vspace{-3mm}
\paragraph{Exercice 7}:  \hfill \textbf{1 points}

Lors d'un set d'un match de tennis, un joueur a réussi $20$ premiers services
sur un total de $28$.
Quel est le pourcentage de réussite au premier service ? Arrondir à l'unité.

\vspace{-3mm}
\paragraph{Exercice 8}:  \hfill \textbf{2 points}

\vspace{-3mm}
\begin{enumerate}
\item Dans un CDI, parmi les $300$ albums jeunesse, $\dfrac{3}{15}$ sont des BD.
Combien y a-t-il de BD dans le CDI ?
\item La catégorie benjamins d'un club de karaté compte $33$ adhérents, dont les deux tiers sont des filles. Combien y a-t-il de filles en catégorie benjamin dans ce club ?
\end{enumerate} 

\vspace{-3mm}
\paragraph{Exercice 9}:  \hfill \textbf{1 points}

Dans un club sportif, $45 \%$ des $260$ adhérents font du tennis.
Quel est le nombre de personnes pratiquant le tennis dans ce club ?

\vspace{-3mm}
\paragraph{Exercice 10}:  \hfill \textbf{1.5 points}

Chloé achète une robe qui coûte $90$\euro . Elle bénéficie d'une réduction de $65 \%$.
\begin{enumerate}
\item Quel est le montant de la réduction ?
\item Combien va-t-elle payer sa robe ?
\end{enumerate}

\vspace{-3mm}
\paragraph{Exercice 11}:  \hfill \textbf{2 points}

\vspace{-3mm}
\begin{enumerate}
\item Lors d'une élection, dans une commune où $480$ votes ont été exprimés, une candidate a obtenu $11,25 \%$ des voix. Calculer le nombre de personnes qui
ont voté pour elle.
\item Pour la même élection, un autre candidat a obtenu $132$ voix. Calculer le pourcentage de votes exprimés pour ce candidat.
\end{enumerate}

\vspace{-3mm}
\paragraph{Exercice 12}:  \hfill \textbf{3 points}

En $2002$, Léo a un salaire de $1 650$\euro . Son patron l'augmente de $8 \%$ en $200$ puis diminue son salaire de $8 \%$ en $2004$. Quel est son salaire en $2004$ ?

\end{document}
